\section{Related Literature}

This study connects three research strands and positions itself at their intersection.

First, the broader technology-productivity literature emphasizes complementarity. Digital technologies produce larger gains when accompanied by organizational adaptation, managerial quality, and human capital investment \cite{brynjolfssonHitt2000,griliches1998}. This perspective is especially relevant for AI, where deployment quality and institutional readiness vary strongly across countries \cite{bresnahanTrajtenberg1995,cominMestieri2018}.

Second, AI-specific macro discussions highlight both upside potential and transitional frictions. Conceptual and policy analyses suggest that AI can raise productivity, but adjustment costs, diffusion lags, and uneven capabilities may delay aggregate growth responses \cite{acemogluRestrepo2018,brynjolfssonRockSyverson2019,aghionJonesJones2017,korinekStiglitz2021,imf2024}. Recent micro-level evidence documents sizable productivity effects from generative AI tools in narrower settings, which helps explain why macro cross-country effects may appear slower and more moderate \cite{noyZhang2023,brynjolfssonLiRaymond2023,eloundouEtAl2023}. As a result, empirical estimates may differ depending on sample period, outcome definition, and the empirical treatment of global shocks.

Third, the panel-identification literature underscores that fixed effects are necessary but not sufficient for causal interpretation. While fixed effects absorb time-invariant heterogeneity, they do not eliminate endogeneity from time-varying confounders, feedback effects, or measurement error \cite{wooldridge2010,wooldridge2021,bai2009,callaway2021}. The implication is methodological: credibility should rely on specification stability, explicit assumptions, and transparent limitations.

Relative to these strands, our paper contributes by combining a reproducible workflow with a robustness-first inferential strategy. Rather than privileging one estimator, we ask whether substantive conclusions remain stable across a structured set of alternatives. This approach is designed to improve cumulative evidence quality in a fast-moving empirical domain where measurement and identification are still evolving.
